% Generated by report_gen.py (ai-ee v1) - do not hand-edit.
\documentclass[11pt,a4paper]{article}
\usepackage[margin=2.2cm]{geometry}
\usepackage{graphicx}
\usepackage{booktabs}
\usepackage{longtable}
\usepackage{pdfpages}
\makeatletter
\@ifundefined{KV@Gin@artifact}{\define@key{Gin}{artifact}[]{}}{}
\makeatother
\usepackage{xcolor}
\usepackage[hidelinks]{hyperref}
\setcounter{tocdepth}{1}
\begin{document}
\sloppy
\begin{center}
{\LARGE\bfseries g0-sense --- Design Document}\\[6pt]
{\large ai-ee v1 pipeline}\\[2pt]
generated 2026-08-27 08:43:41
\end{center}
\tableofcontents

\section{Board and Run Metadata}
{\small
\begin{longtable}{lp{11cm}}
\toprule
\textbf{Field} & \textbf{Value} \\
\midrule
\endhead
board & g0-sense \\
workspace & \texttt{boards/g0-sense} \\
phase & P2 \\
gate status & no gates recorded yet \\
state created & 2026-08-27T06:35:38 \\
state updated & 2026-08-27T08:42:51 \\
\bottomrule
\end{longtable}
}
\section{Overview}
\subsection*{Brief: g0-sense (unattended container run, 2026-08-27)}
A small 2-layer USB-C powered temperature/humidity sensor node - a product-scope board (no mode token: design normally), meant to be sent to JLCPCB as-is.

\begin{itemize}
\item Power: USB-C receptacle used for POWER ONLY (5 V VBUS; 5.1 k CC pull-downs so
any USB-C source or a legacy A-to-C cable enables VBUS; D+/D- unconnected).
Sensible input protection for a USB-powered device (VBUS TVS, resettable
fuse or equivalent, reverse/inrush as the architect sees fit).
\item 3.3 V rail from a fixed LDO, \textgreater{}= 300 mA, JLC Basic/Preferred part if one fits
(AMS1117-3.3 class is fine), with proper input/output capacitors.
\item MCU: STM32G030F6P6 (TSSOP-20) on the internal oscillator - no crystal.
Decoupling per the datasheet, NRST with a push button, BOOT0 handled.
\item Sensor: Sensirion SHT40 (or the SHT4x family member in JLC stock) on I2C with
pull-ups; give it a thermal isolation slot/cutout or at least keep it away
from the LDO and MCU heat.
\item Interfaces: 4-pin SWD header (0.1 in), 4-pin UART header (0.1 in: GND, 3V3,
TX, RX) for readout/logging, one Qwiic/STEMMA-QT connector (JST SH 1.0 mm
4-pin, GND/3V3/SDA/SCL) sharing the sensor I2C bus.
\item One user LED (GPIO, series resistor) and one power LED.
\item 2 layers, 1.6 mm FR-4, HASL, green; JLC economy PCBA for the SMD parts,
prototype qty 5. Prefer JLC Basic parts; Extended parts only where the brief
names the part (MCU, sensor, connectors).
\item Size: whatever an honest layout wants, roughly 35 x 25 mm or smaller; four
M2 mounting holes if they fit without hurting the layout.
\item No battery, no mains, no USB data, no RF.
\end{itemize}
\subsection*{Stackup}
Chosen stackup: \texttt{JLC2313\_1.6} (reference/stackups.yaml, \texttt{available: true})
\section{Requirements}
\subsection*{requirements - g0-sense}
P0 requirements distilled from \texttt{brief/brief.md} (2026-08-27, unattended container run).

\subsubsection*{1. Function}
Mode declaration: the brief carries NO mode token, so mode = none - design normally. Scope = product-scope per the brief's own words ("a product-scope board ... meant to be sent to JLCPCB as-is"): protection, filtering, connectors, thermal and enclosure-fit are all in scope, and their absence is a reviewer finding. Binding: the stated \textasciitilde{}35 x 25 mm size is a SOFT preference, not a hard cap (see section 5). Stage under study: none.

The board is a small 2-layer USB-C powered temperature/humidity sensor node. A USB-C receptacle supplies 5 V (power only, no data); a fixed 3.3 V LDO powers an STM32G030F6P6 MCU running on its internal oscillator, which reads a Sensirion SHT4x sensor over I2C. The same I2C bus is exposed on a Qwiic/STEMMA-QT connector; readout/logging is over a 4-pin UART header; programming/debug is over a 4-pin SWD header. One user LED and one power LED. Prototype quantity 5, JLC economy PCBA, meant to be ordered as-is.

\subsubsection*{2. Interfaces}
\begin{itemize}
\item USB-C receptacle, POWER ONLY: 5 V VBUS input; 5.1 k pull-downs on CC1 and
CC2 (so any USB-C source or a legacy A-to-C cable enables VBUS); D+/D-
unconnected; no USB data.
\item SWD header: 4-pin, 0.1 in pitch. ASSUMED: pins GND, 3V3, SWDIO, SWCLK
(exact order per architect; brief states only "4-pin SWD header (0.1 in)").
\item UART header: 4-pin, 0.1 in pitch: GND, 3V3, TX, RX - for readout/logging.
ASSUMED: TX/RX are named from the MCU's perspective (TX = MCU transmit).
\item Qwiic/STEMMA-QT connector: JST SH 1.0 mm 4-pin, GND/3V3/SDA/SCL, sharing
the sensor I2C bus. I2C pull-ups required on the board (values per
architect).
\item User LED: one, driven from a GPIO through a series resistor.
\item Power LED: one (which rail it indicates - VBUS or 3V3 - is per architect;
the brief does not say).
\item NRST push button (momentary reset).
\item BOOT0: handled (method - strap, option bits, or jumper - per architect; the
brief requires only that it is handled).
\item Four M2 mounting holes, CONDITIONAL: only if they fit without hurting the
layout (see section 5).
\end{itemize}
\subsubsection*{3. Power}
\begin{itemize}
\item Input: USB-C VBUS, 5 V nominal, device draws \textless{} 1 A total. No battery, no
charging, no mains.
\item VBUS input protection required (product scope, brief-stated): TVS,
resettable fuse or equivalent; reverse/inrush handling as the architect
sees fit - the brief delegates the exact topology.
\item 3.3 V rail: fixed LDO, \textgreater{}= 300 mA. JLC Basic/Preferred part preferred;
AMS1117-3.3 class acceptable. Proper input/output capacitors per the chosen
part's datasheet.
\item Rail budget (GUESSES - to be replaced by the architect's numbers):
\begin{itemize}
\item STM32G030F6P6 on internal oscillator: \textasciitilde{}5-15 mA (guess)
\item SHT4x: \textless{} 1 mA average, single-digit mA peaks during measurement (guess)
\item LEDs: \textasciitilde{}2-5 mA total (guess)
\item Qwiic downstream devices: not budgeted by the brief - see open
question 2 (recommended default: reserve 100 mA).
\item On-board total is tens of mA; the brief's \textgreater{}= 300 mA LDO floor leaves the
headroom for Qwiic loads.
\end{itemize}
\end{itemize}
\subsubsection*{4. Environment}
Not stated in the brief. ASSUMED: indoor, room-ambient use (roughly 0-40 C), bare board, no ingress or vibration requirements (see open question 1). The one stated thermal/environmental requirement is layout-level: the SHT4x gets a thermal isolation slot/cutout, or at minimum is kept away from LDO and MCU heat, so self-heating does not corrupt the measurement.

\subsubsection*{5. Size \& mounting}
\begin{itemize}
\item Outline: "whatever an honest layout wants, roughly 35 x 25 mm or smaller".
This is a SOFT preference - no HARD cap. It must NOT bind at P5
board\_init; the layout earns the final outline, and \textasciitilde{}35 x 25 mm is a
target to aim for, not a box to squeeze into (see open question 5).
\item Mounting: four M2 mounting holes, CONDITIONAL - "if they fit without
hurting the layout". The brief itself allows dropping or reducing them in
favor of layout honesty; record that decision if taken.
\item Height: no limit stated.
\end{itemize}
\subsubsection*{6. Quantity \& budget}
\begin{itemize}
\item Quantity: 5 boards, prototype run, assembled.
\item Unit cost: no target stated (see open question 3). Cost posture from the
brief: JLC economy PCBA; prefer JLC Basic parts; Extended parts only where
the brief names the part (MCU, sensor, connectors).
\end{itemize}
\subsubsection*{7. Assembly}
\begin{itemize}
\item JLC economy PCBA for the SMD parts.
\item ASSUMED: single-sided SMD assembly (economy PCBA norm, and this part count
does not need the second side).
\item Through-hole parts: economy PCBA assembles SMD only, so the two 0.1 in
headers are ASSUMED to ship unpopulated for owner hand-soldering (see open
question 4).
\item ASSUMED: NRST push button and the USB-C / JST SH connectors are SMD,
JLC-assemblable variants where stock allows, so they ride the PCBA run.
\item Board: 2 layers, 1.6 mm FR-4, HASL, green (brief-stated).
\end{itemize}
\subsubsection*{8. Compliance/safety flags}
Each flag checked against the brief; none applies:

\begin{itemize}
\item Mains voltage: NOT APPLICABLE - brief states "no mains"; input is USB 5 V.
\item Battery: NOT APPLICABLE - brief states "no battery"; no charging circuitry
anywhere in scope.
\item Motors: NOT APPLICABLE - no motor or actuator in the brief's scope.
\item \textgreater{}30 V: NOT APPLICABLE - the highest voltage on the board is USB VBUS at
5 V nominal.
\item High current (\textgreater{}3 A): NOT APPLICABLE - the brief describes a USB 5 V,
\textless{} 1 A device.
\item RF transmit: NOT APPLICABLE - brief states "no RF"; no radio, and no USB
data either ("D+/D- unconnected").
\end{itemize}
\subsubsection*{9. Open questions}
Closed-form, each with a recommended default. Unattended run: the orchestrator answers these on the owner's behalf and records the answers.

1. Operating environment: is indoor, room-ambient use (about 0-40 C, bare board, no ingress/vibration requirements) correct? Recommended default: YES. 2. Qwiic 3.3 V budget: how much current may external Qwiic/STEMMA-QT devices draw from the 3.3 V rail? Recommended default: reserve 100 mA (fits under the \textgreater{}= 300 mA LDO floor together with the on-board load). 3. Unit cost: is there a hard per-board cost cap? Recommended default: NO hard cap - minimize via economy PCBA and Basic parts as the brief already directs. 4. Header population: should the two 0.1 in headers (SWD, UART) ship unpopulated for hand-soldering, given economy PCBA assembles SMD only? Recommended default: YES - footprints and holes present, headers not populated at assembly. 5. Size confirmation: is \textasciitilde{}35 x 25 mm truly a soft preference (the layout may exceed it if honesty demands), not an enclosure-driven hard cap? Recommended default: SOFT preference, per the brief's own wording ("whatever an honest layout wants").

\section{Architecture}
\subsection*{Decisions of record (state.json)}
{\small
\begin{longtable}{lp{5.2cm}p{7.2cm}}
\toprule
\textbf{Phase} & \textbf{What} & \textbf{Why} \\
\midrule
\endhead
P0 & unattended default: environment = indoor bench/room ambient, 0-40 C operating, bare board, no ingress or vibration rating & Q1 of requirements section 9; brief names no environment. Most conservative common-practice choice for a USB-powered indoor sensor node; parts will still be picked at industrial -40..85 C grade where that is the stocked option, so the assumption cannot make the board wrong. \\
P0 & unattended default: Qwiic/STEMMA-QT downstream budget = 100 mA reserved on the 3V3 rail & Q2; brief does not budget it. 100 mA covers typical Qwiic sensor/display loads and still fits inside the brief's \textgreater{}= 300 mA LDO floor with margin for the on-board tens of mA, so it sizes the LDO and its thermals honestly without inflating the part. \\
P0 & unattended default: no hard per-board cost cap; minimize cost via JLC economy PCBA and Basic parts as the brief already directs & Q3; brief states no budget. Cost is already bounded by the brief's own Basic-parts-preferred rule; inventing a cap could force a worse engineering choice. \\
P0 & unattended default: the two 0.1 in headers (SWD, UART) ship UNPOPULATED - plated through-holes only, hand-solderable; only SMD parts go to JLC economy PCBA & Q4; brief says 'JLC economy PCBA for the SMD parts'. JLC economy assembly is SMT-only, so THT headers are BOM 'do not populate / customer supplied'. Standard practice for dev/sensor boards and it keeps the assembly quote in economy tier. \\
P0 & unattended default: \textasciitilde{}35 x 25 mm is a SOFT preference, not an enclosure-driven hard cap; no HARD outline cap binds at P5 & Q5; the brief says 'whatever an honest layout wants, roughly 35 x 25 mm or smaller' and names no enclosure. Confirms requirements.md section 5; placement must not optimize to fit an unearned number (the bb-buck failure mode). \\
P1 & unattended default: LDO = AMS1117-3.3 class in SOT-223 (JLC Basic, C6186); governing thermal design point is the brief's literal \textgreater{}=300 mA rated case (0.51 W), not the \textasciitilde{}150 mA realistic load & research/ldo.md escalated the SOT-223-vs-SOT-23-5 tradeoff; research/power.md resolved it with numbers: AMS1117 Tj 71-86 C at 40 C ambient and \textless{}=115 C even in the 500 mA abuse case, while AP2112K-class SOT-23-5 (theta-JA 184 C/W) hits Tj 134 C at the rated case with no copper fix available. The Qwiic port is an unfenced 3V3 output, so the rated case is reachable in the field. Basic part, 1.1M stock, \$0.20. \\
P1 & unattended default: a sub-6 V TVS clamping voltage is NOT a requirement; SMAJ5.0A/SMF5.0A-class (standoff 5.0-5.5 V, Vc \textasciitilde{}9.2 V at Ipp) is accepted & research/protection.md flagged that no passive TVS clamps under 6 V. Nothing on VBUS is 6 V-rated: the only VBUS consumers are the LDO (AMS1117 Vin abs max 15 V), the input caps (16 V+ parts) and the PTC (16 V). Everything at 3.3 V sits behind the LDO. A \textasciitilde{}9.2 V clamp during a surge is therefore inside every VBUS part's rating - the original \textless{}6 V framing was an over-tight derived spec, not a real one. \\
P1 & unattended default: NO series reverse-polarity element (no Schottky, no P-FET ideal diode) on VBUS & Two independent reasons, both recorded in research: (a) a compliant USB-C cable cannot present reverse polarity (protection.md), and (b) research/power.md computed the end-to-end dropout budget and found a series Schottky's 125-190 mV would consume the entire AMS1117 headroom at the 4.75 V source corner (+0.11 V worst-corner margin as designed). The unidirectional TVS already crowbars a reverse fault. Adding it would make the board worse, not safer. \\
P1 & unattended default: USB-C receptacle = 16-pin power-only part whose THT shield/mounting legs are accepted; JLC assembles it in the Economic PCBA tier & research/connectors.md correctly flagged that every 16P part in this class has THT shield legs and asked whether economy PCBA can place them. Verified directly on the JLCPCB part page for C165948 (TYPE-C-31-M-12): Assembly Process = SMT, Mount Type = Surface Mount Right Angle, and the part is listed as supported by BOTH Economic and Standard assembly. So the legs are not an assembly blocker. Part is JLC Extended - permitted, the brief names the connector. \\
P1 & unattended default: SHT4x on-chip heater (75 mA peak) IS budgeted in the rail tree at \textless{}=10\% duty & research/power.md's call, accepted: including it costs nothing (peak +3V3 still 185 mA, 240 mA with 30\% headroom, under the 300 mA LDO floor and at 39\% of the 500 mA CC-blind USB entitlement), while omitting it would turn a documented firmware feature into a hardware ECO later. \\
P1 & unattended default: NRST tactile button - prefer the JLC Basic 5.1x5.1 mm part over the brief's smaller 3x2.5/4.2x3.2 mm size class if the layout has room; size class is a preference, Basic status and stock depth are the tiebreak & research/connectors.md Q2: the only Basic button is 5.1x5.1 mm (C318884, 962k stock, \$0.02); the closest-size alternative carries a documented lib\_pull symbol gotcha in this repo's own LEARNINGS. The brief never fixed the button size, and the board size is soft, so paying \textasciitilde{}2.9 mm of edge for a Basic, gotcha-free part is the conservative trade. Architect may overrule at P2 if placement suffers. \\
P1 & unattended default: LED colours - user LED green, power LED red; both 0603, series resistors sized from the measured Vf band per colour & research/connectors.md Q3 left the assignment open. Red is the conventional power indicator and is the JLC Basic option with the low Vf band (1.8-2.4 V), which gives the largest resistor and the lowest current draw for an always-on LED. Green for the user LED is the conventional 'activity' colour; it is Extended in the KT-0603 family, which is acceptable at one extra unique part, and its 2.6-3.1 V Vf is fine from 3.3 V at a few mA. \\
P2 & research task block-mcu-1 closed over 1 unverified draft record(s): mcu-boot-strap-only-when-option-bytes-consult-pin & --accept-drafts: the slot ships without this knowledge \\
P2 & unattended default: BOOT0 gets a populated 10k pull-down on PA14 (pin 19, shared with SWCLK) - the board does NOT rely on the unverified factory-option-byte claim & The research record mcu-boot-strap-only-when-option-bytes-consult-pin was REFUTED by a fresh second reader: its load-bearing claim (factory 0xDFFFE1AA / nBOOT\_SEL=1 makes the BOOT0 pin ignored) rests solely on a community.st.com forum page, and RM0454/AN2606 could not be acquired - st.com times out from this container (verified twice by curl, once by an agent). Rather than design on unverified knowledge, the board is made correct under EITHER option-byte state: a 10k pull-down guarantees BOOT0=0 at reset so the part boots main flash even if nBOOT\_SEL=0. Cost is one 0402 resistor; 330 uA load on SWCLK is negligible to any debugger's push-pull driver, and the datasheet-confirmed internal pull-down on PA14 in debug mode already pulls the same way. This closes the gap without the source. \\
P2 & unattended default: I2C pull-ups are 1.5 kOhm (not the 2.2 kOhm the P1 fragment proposed) & The draft record sht4x-i2c-pullup-bus-cap was refuted for recommending 2.2k over its own declared 150-200 pF bus budget: UM10204 Eq 1 at tr=300 ns gives Rp(max) = 354/Cb{[}pF{]} kOhm, i.e. 1.77k at 200 pF, so 2.2k needs Cb \textless{}= 161 pF and misses the Fast-mode rise time on a cabled Qwiic bus. 1.5k holds the ceiling to 236 pF and clears both floors (967 Ohm from UM10204 Table 10 at 3.3 V/3 mA, 390 Ohm from the SHT4x datasheet). Value independently recomputed by a second reader. architecture/constraints.json and the P4 schematic must use 1.5k. \\
\bottomrule
\end{longtable}
}
\subsection*{Sheet plan}
\subsection*{sheets - g0-sense hierarchical plan (P2)}
Small board: root + TWO child sheets. One sheet would also fit, but the power/main split gives P6 its natural placement groups (power entry corridor vs logic/sensor), keeps the Type-C + protection cluster reviewable in one frame, and costs nothing. Refdes are unique ACROSS sheets by disjoint ranges; every range includes a \#PWR base.

\subsubsection*{Root: \texttt{g0-sense.kicad\_sch}}
Holds the two sheet symbols only. Global power nets (power symbols, bare netlist names): \texttt{VBUS}, \texttt{+5V}, \texttt{+3V3}, \texttt{GND}. Power symbols make these names bare in the final netlist; all constraints.json net references use the bare forms.

\subsubsection*{Sheet 1: \texttt{power.kicad\_sch} (blocks B1 usb-input+protection, B2 ldo-3v3)}
\begin{itemize}
\item Contents: J1 USB-C receptacle; R1, R2 (5.1k Rd); D1 (TVS); F1 (PTC);
C1 (100 nF VBUS); U1 (AMS1117-3.3); C2 (10 uF X5R, +5V at VIN);
C3 (22 uF tantalum, +3V3); D2 (power LED red) + R3 (620 Ohm).
\item Interface nets (hier pins -\textgreater{} placement groups at P6): \texttt{+3V3} (out),
\texttt{GND}. VBUS and +5V stay sheet-local (they leave via power symbols, so
their netlist names are still bare \texttt{VBUS} / \texttt{+5V}).
\item Sheet-local labels: \texttt{CC1}, \texttt{CC2} (final names \texttt{/power/CC1}, \texttt{/power/CC2}).
\item Refdes ranges: J1; U1; F1; D1-D9; R1-R9; C1-C9; SW none; \#PWR pwr\_base=100
(\#PWR0100-\#PWR0199).
\end{itemize}
\subsubsection*{Sheet 2: \texttt{main.kicad\_sch} (blocks B3 mcu, B4 sht4x + I2C/Qwiic, headers)}
\begin{itemize}
\item Contents: U2 (STM32G030F6P6); U3 (SHT40-AD1B class); C10 (100 nF U2),
C11 (4.7 uF U2), C12 (100 nF NRST), C13 (100 nF U3); R10, R11 (1.5k I2C
pull-ups), R12 (220 Ohm user LED), R13 (10k BOOT0 pull-down); D10 (user LED green); SW1 (NRST
button); J2 (Qwiic JST SH); J3 (SWD 1x4 DNP); J4 (UART 1x4 DNP);
H1-H4 (M2 mounting holes, conditional).
\item Interface nets: \texttt{+3V3} (in), \texttt{GND}.
\item Sheet-local labels (final names \texttt{/main/\textless{}label\textgreater{}}): \texttt{SDA}, \texttt{SCL}, \texttt{SWDIO},
\texttt{SWCLK}, \texttt{UART\_TX}, \texttt{UART\_RX}, \texttt{NRST}, \texttt{LED\_USER}.
\item Refdes ranges: U2-U9; J2-J9; D10-D19; R10-R19; C10-C19; SW1-SW9; H1-H9;
\#PWR pwr\_base=200 (\#PWR0200-\#PWR0299).
\end{itemize}
\subsubsection*{Interface-net contract (P4 must produce exactly these)}
\begin{quote}\small\ttfamily\raggedright
\textbar{} Net (netlist form) \textbar{} Class \textbar{} From \textbar{} To \textbar{}\\
\textbar{}---\textbar{}---\textbar{}---\textbar{}---\textbar{}\\
\textbar{} VBUS \textbar{} power (bare) \textbar{} J1 \textbar{} D1, F1, C1 \textbar{}\\
\textbar{} +5V \textbar{} power (bare) \textbar{} F1 \textbar{} U1 VIN, C2 \textbar{}\\
\textbar{} +3V3 \textbar{} power (bare) \textbar{} U1 VOUT \textbar{} C3, D2/R3, U2, U3, R10, R11, J2, J3.2, J4.2 \textbar{}\\
\textbar{} GND \textbar{} power (bare) \textbar{} all \textbar{} all + J1 shield \textbar{}\\
\textbar{} /power/CC1, /power/CC2 \textbar{} signal \textbar{} J1 CC pins \textbar{} R1, R2 \textbar{}\\
\textbar{} /main/SDA, /main/SCL \textbar{} signal \textbar{} U2 PA10/PA9 \textbar{} U3, J2, R10/R11 \textbar{}\\
\textbar{} /main/SWDIO, /main/SWCLK \textbar{} signal \textbar{} U2 PA13/PA14 \textbar{} J3.3/J3.4; SWCLK also R13 -\textgreater{} GND \textbar{}\\
\textbar{} /main/UART\_TX, /main/UART\_RX \textbar{} signal \textbar{} U2 PA2/PA3 \textbar{} J4.3/J4.4 \textbar{}\\
\textbar{} /main/NRST \textbar{} signal \textbar{} U2 pin 6 \textbar{} C12, SW1 \textbar{}\\
\textbar{} /main/LED\_USER \textbar{} signal \textbar{} U2 PA5 \textbar{} R12 -\textgreater{} D10 \textbar{}
\end{quote}
Header pinouts (canonical): J3 SWD = 1:GND, 2:3V3, 3:SWDIO, 4:SWCLK. J4 UART = 1:GND, 2:3V3, 3:TX (MCU transmit), 4:RX (MCU receive). Pin 1 marked on silk on both; no vendor standard exists for a 4-pin SWD header, so the silk labels are the contract.

Full architecture narrative (not inlined; see the artifact index): \texttt{architecture/blocks.md}, \texttt{architecture/decisions.md}, \texttt{architecture/power\_tree.md}, \texttt{architecture/stackup.md}.
\section{Schematic}
\emph{Pending --- produced at P4.}
\section{Layout}
\emph{Pending --- produced at P6.}
\section{Verification}
\emph{Pending --- produced at P8.}
\section{DFM and Fabrication}
\emph{Pending --- produced at P9.}
\section{Run Record (Appendix)}
\subsection*{Phase digests}
\paragraph*{P0}
\begin{quote}\small\ttfamily\raggedright
\# P0 Intake digest - g0-sense (2026-08-27)\\
~\\
- Artifact: `requirements.md` (9 sections), `check\_requirements.py` PASS (0 violations).\\
- Mode: NO mode token in the brief -\textgreater{} mode = none, design normally. Scope is\\
  product per the brief's own words, so absent protection/filtering/connectors/\\
  thermal IS a reviewer finding downstream.\\
- Geometry: \textasciitilde{}35 x 25 mm recorded SOFT (no HARD cap) - must not bind at P5\\
  board\_init; four M2 holes conditional on the layout.\\
- Safety (section 8): mains / battery / motors / \textgreater{}30 V / \textgreater{}3 A / RF transmit all\\
  NOT APPLICABLE with brief evidence (USB-C 5 V power-only, \textless{}1 A). No safety\\
  unknown to escalate; nothing PROVISIONAL on safety grounds.\\
- 5 OPEN questions answered as unattended defaults and recorded as decisions:\\
  environment indoor 0-40 C; Qwiic budget 100 mA; no hard cost cap; 0.1 in\\
  headers ship unpopulated (economy PCBA is SMT-only); size stays soft.\\
- Spawn: requirements-analyst @ fable/high (session effort xhigh).\\
- Next: P1 research roster (power architect, component scout, interface spec).
\end{quote}
\paragraph*{P1}
\begin{quote}\small\ttfamily\raggedright
\# P1 Research digest - g0-sense (2026-08-27)\\
~\\
Roster (8 agents): 4x research-component-scout (ldo, protection, connectors,\\
mcu-sensor), research-interface-spec (usbc-power-sink), 2x\\
research-reference-design (sht4x, stm32g030-minimal), research-power-architect.\\
~\\
Load-bearing findings\\
- USB-C CC-blind sink (fixed 5.1k Rd on CC1 AND CC2, independent - never one\\
  shared resistor) is entitled to DEFAULT USB power only: budget \textless{}= 500 mA.\\
  Sink capacitance at the receptacle \textless{}= 10 uF (TC2.0 Table 4-3).\\
- Rail tree: VBUS -\textgreater{} TVS -\textgreater{} PTC -\textgreater{} +5V -\textgreater{} LDO -\textgreater{} +3V3. +3V3 peak 185 mA\\
  (incl. 100 mA Qwiic reserve + 75 mA SHT4x heater), 240 mA with 30\% headroom.\\
  VBUS worst draw 196 mA = 39\% of entitlement.\\
- LDO: AMS1117-3.3 SOT-223 (C6186, Basic) beats SOT-23-5 low-Iq parts on the\\
  only axis that decides: Tj 71-86 C vs 134 C at 40 C ambient, 0.51 W case.\\
  Needs a 22 uF TANTALUM-class output cap (not plain ceramic) - carry to P3/P4.\\
  Dropout closes at the 4.75 V corner with +0.11 V worst-case margin.\\
- STM32G030F6P6 TSSOP-20: VDD/VDDA bonded to pin 4, VSS/VSSA pin 5, no separate\\
  VREF+. BOOT0 needs NO strap - factory nBOOT\_SEL=1/nBOOT0=1 always boots main\\
  flash, and BOOT0 shares pin 19 with SWCLK anyway. NRST needs only 100 nF + the\\
  button (internal pull-up 40k typ). Pin map: SWD PA13/PA14, USART2 PA2/PA3,\\
  I2C1 PA9/PA10 (5 V tolerant, Fm+), LED PA5. 9 of 20 pins used, no conflicts.\\
- SHT4x: 100 nF at VDD, 2.2k I2C pull-ups (sized for a 0.5 m Qwiic leg), copper\\
  conduction is the dominant self-heating path -\textgreater{} distance + necked traces +\\
  no pour under the part + a milled slot. NO BOARD WASH (Sensirion forbids it)\\
  -\textgreater{} must appear in the fab/assembly notes. Fixed I2C address per variant.\\
- MCU + sensor are both JLC EXTENDED (no Basic exists in either family):\\
  STM32G030F6P6 C724040 \$1.37, SHT40-AD1B-R2 C2909890 \$1.90 (addr 0x44).\\
~\\
7 OPEN questions answered and recorded as decisions (LDO design point + part;\\
TVS clamp ceiling dropped as an over-tight derived spec; no reverse-polarity\\
element; USB-C THT legs verified OK in JLC Economic tier; heater budgeted;\\
button size vs Basic; LED colours). Gates: none due at P1.
\end{quote}
\paragraph*{P2}
\emph{(no digest recorded)}
\subsection*{Gate attempt history}
\emph{(no entries)}
\subsection*{Issues}
\emph{(no entries)}
\subsection*{Human checkpoints}
{\small
\begin{longtable}{lllp{2.6cm}p{6.4cm}}
\toprule
\textbf{Id} & \textbf{Phase} & \textbf{Status} & \textbf{When} & \textbf{Note} \\
\midrule
\endhead
1 & P2 & not reached &  &  \\
2 & P4 & not reached &  &  \\
3 & P6 & not reached &  &  \\
4 & P8 & not reached &  &  \\
5 & P10 & not reached &  &  \\
\bottomrule
\end{longtable}
}
\section{Artifact Index}
{\small
\begin{longtable}{p{9cm}rp{4cm}}
\toprule
\textbf{File} & \textbf{Size} & \textbf{Note} \\
\midrule
\endhead
\texttt{state.json} & 19,847 B &  \\
\texttt{brief/brief.md} & 1,634 B &  \\
\texttt{requirements.md} & 6,653 B &  \\
\texttt{architecture/blocks.md} & 11,799 B &  \\
\texttt{architecture/constraints.json} & 8,406 B &  \\
\texttt{architecture/decisions.md} & 7,408 B &  \\
\texttt{architecture/power\_tree.md} & 4,040 B &  \\
\texttt{architecture/sheets.md} & 2,964 B &  \\
\texttt{architecture/stackup.md} & 2,679 B &  \\
\bottomrule
\end{longtable}
}
\end{document}
